problem:  
Let \((M^{2n+1}, T^{1,0}M, \theta)\) be a closed strictly pseudoconvex pseudohermitian manifold with Tanaka–Webster connection \(\nabla^{\mathrm{TW}}\), Levi form \(h_{\alpha\bar\beta}\), and torsion tensor \(A_{\alpha\beta}\).  
Let \(\Delta_b\) denote the sub-Laplacian, and for each normalized eigenfunction \(f_j\) satisfying  
\[
\Delta_b f_j = -\lambda_j f_j, \qquad \|f_j\|_{L^2(M,\theta)} = 1,
\]
consider its microlocal energy distribution (Wigner distribution) \(\mu_j\) on the horizontal (sub-Riemannian) cotangent bundle \(H^*M := \{(x,\xi)\in T^*M : \xi(\,T\,)=0\}\), where \(T\) is the Reeb vector field.  
The principal symbol of \(\Delta_b\) is \(p(x,\xi) = |\xi|_{h^{-1}}^2\) on \(H^*M\), generating the **horizontal geodesic (Hamiltonian) flow** \(\varphi_t\) on \(p^{-1}(1)\).

Any weak-* limit point \(\mu\) of the sequence \((\mu_j)\) as \(\lambda_j\to\infty\) is called a *CR semiclassical measure* or *quantum limit*.  
Such measures are expected (by sub-Riemannian analogues of Egorov and propagation laws) to be invariant under \(\varphi_t\).

**Problem:**  
Determine whether the Tanaka–Webster torsion \(A_{\alpha\beta}\) leaves a detectable imprint on the possible CR quantum limits.  
Precisely:
- Does every semiclassical measure \(\mu\) remain supported on the set of *torsion-null directions*, i.e. those \((x,\xi)\in p^{-1}(1)\) for which the torsion-contracted quadratic form \(A_{\alpha\beta}(x)\, h^{\alpha\bar\gamma}h^{\beta\bar\delta}\,\xi_{\bar\gamma}\xi_{\bar\delta}=0\)?  
- Or can there exist nontrivial pseudohermitian manifolds (with \(A\neq 0\)) admitting sequences of sub-Laplacian eigenfunctions whose semiclassical measures \(\mu\) give positive weight to *torsion-active directions* (where that contraction is nonzero)?  
Equivalently, does torsion appear only in subprincipal corrections invisible to quantum limits, or can it alter the invariant dynamics of these measures?  

A rigorous result in either direction would clarify whether pseudohermitian torsion affects the high-frequency limit of CR sub-Laplacian eigenfunctions at the microlocal level.

Why is it a "good" problem:  
This problem asks a deep and natural question at the crossroads of **CR geometry, sub-Riemannian analysis, and semiclassical microlocal theory**: does pseudohermitian torsion—one of the key geometric features distinguishing CR manifolds from contact Riemannian ones—impact the asymptotic microlocal distribution of sub-Laplacian eigenfunctions?  
It translates a geometric structure (torsion) into a precise analytic question about invariant measures under Hamiltonian flows associated with \(\Delta_b\).  
The formulation is conceptually simple yet analytically profound: it demands understanding of the subprincipal symbol of \(\Delta_b\), the propagation of microlocal mass in a hypoelliptic setting, and the potential emergence of torsion-dependent invariant distributions.  
Its resolution could extend **quantum ergodicity** and **semiclassical analysis** to the CR context, revealing whether torsion plays a measurable spectral role or remains invisible in the high-energy limit—an open and fundamentally geometric problem with broad implications.